% ============================================================================
%  অধ্যায় ২ — বীজগণিত ১
%  COMPILE WITH XeLaTeX  (twice, so the point total resolves)
%  Overleaf: Menu -> Compiler -> XeLaTeX
% ============================================================================
\documentclass[11pt]{article}
\usepackage{bnhandout}          % add [nosol] for a student edition

\hauthor[https://jugantordey.github.io/]{Jugantor Dey}
\hseries{\textit{Mathematics for the Physics Olympiad}}

\begin{document}

\handouttitle{Chapter 2 : Algebra I}

\noindent
এই অধ্যায়ের একমাত্র পূর্বশর্ত অধ্যায় ১। সেখান থেকে দুটো জিনিস এখানে
বারবার লাগবে --- একটিমাত্র counterexample দিয়ে যে কোনো দাবি ভেঙে ফেলা যায়, আর
$\Rightarrow$ ও $\Leftrightarrow$ চিহ্ন দুটোর পার্থক্য। এর বাইরে পাটিগণিত আর
সাধারণ ভগ্নাংশের হিসাব জানা থাকলেই যথেষ্ট।

\noindent বীজগণিত মানে সংখ্যার বদলে অক্ষর বসিয়ে অঙ্ক করা নয়। বীজগণিত হলো একটা কথা একবারেই
সব সংখ্যার জন্য বলে ফেলার উপায়। ``দুই আর তিনের যোগফল তিন আর দুইয়ের যোগফলের
সমান'' --- এটা একটা তথ্য। ``$a + b = b + a$'' --- এটা অসংখ্য তথ্য একসাথে।
এই পার্থক্যটাই এই অধ্যায়ের আসল বিষয়। এখানে মোট \totalpoints\ পয়েন্ট আছে।

\section{চলক ও ধ্রুবক}

\begin{idea}[চলক একটি ফাঁকা ঘর]
একটি চলক (variable) কোনো লুকিয়ে থাকা সংখ্যা নয়, যেটা খুঁজে বের করতে হবে। চলক
হলো একটা ফাঁকা ঘর; একটা জায়গা, যেখানে সংখ্যা বসানো যায়।

$3x + 5$ লেখাটা নিজে কোনো সংখ্যা নয়। এটা একটা নির্দেশনা: ``যে সংখ্যাই বসাও, তাকে
$3$ দিয়ে গুণ করে $5$ যোগ করো''। ঘরে অর্থাৎ $x$ -এর জায়গায় সংখ্যা বসানোর পরেই কেবল এটা সংখ্যা হয়।

(উল্লেখ্য, বীজগণিতে একটি সংখ্যার সাথে বা একটি চলকের সাথে অন্য এক বা একাধিক চলকের গুণন বোঝাতে তাদেরকে পাশাপাশি লেখা হয়। যেমন $3 \times x=3x, 5 \times b=5b, a \times x=ax, p \times q \times r=pqr$ ইত্যাদি।)
\end{idea}

\noindent
একই অক্ষর প্রসঙ্গভেদে তিন রকম কাজ করে, আর এই তিনটা গুলিয়ে ফেলা শুরুর দিকের
সবচেয়ে সাধারণ বিভ্রান্তি।

\begin{enumerate}
  \item \textbf{সব সংখ্যার প্রতিনিধি :} $a + b = b + a$ লেখায় $a$ আর $b$ যেকোনো
        সংখ্যা হতে পারে। এখানে কিছু ``বের করার'' বা নির্ণয় করার নেই।
  \item \textbf{অজানা রাশি :} $3x + 5 = 17$ -এতে $x$ যেকোনো সংখ্যা হতে পারে না;
        কেবল যেসব সংখ্যা গাণিতিক বাক্যটা সত্য করে, সেগুলোই চলবে। এখানে বের করার আছে।
  \item \textbf{ধ্রুবক হিসেবে ব্যবহৃত অক্ষর :} $ax + b = 0$ লেখায় $x$ অজানা,
        কিন্তু $a$ ও $b$ নির্দিষ্ট সংখ্যা, আমরা শুধু জানি না এরা কোন সংখ্যা।
        এদের বলে পরামিতি (parameter)। 
        
\end{enumerate}
 কে variable আর কে parameter, তা কোন অক্ষর ব্যবহার করা হচ্ছে তার উপর নির্ভর করে না, নির্ভর করছে প্রসঙ্গের উপর। যেমন, (ক)-তে $a, b$ variable হলেও (গ)-তে $a, b$ parameter। এব্যাপারটা প্রথম দিকে বুঝতে একটু কষ্ট হলেও সময়ের সাথে সাথে বুঝে ফেলতে পারবে।
        
\begin{remark}[সাবস্ক্রিপ্ট]
একই ধরনের অনেকগুলো রাশি লাগলে নতুন নতুন অক্ষর না খুঁজে সাবস্ক্রিপ্ট ব্যবহার করাই
রীতি: $x_0, x_1, x_2, x_3, \dots$। এখানে $x_1$ পুরোটাই একটিমাত্র variable; $x$ গুণ $1$ নয়।
\end{remark}

\begin{example}
একটি আয়তক্ষেত্রের দৈর্ঘ্য তার প্রস্থের দ্বিগুণের চেয়ে $3$ একক বেশি। প্রস্থকে
$b$ ধরে পরিসীমা $P$-কে $b$-এর মাধ্যমে প্রকাশ করো। এরপর (ক) $b = 5$ হলে $P$
নির্ণয় করো, (খ) $P = 48$ হলে $b$ নির্ণয় করো।
\end{example}

\begin{solbox}
দৈর্ঘ্য $= 2b + 3$। সুতরাং
\[
  P = 2\big(b + (2b+3)\big) = 2(3b + 3) = 6b + 6 .
\]

(ক) $b = 5$ বসিয়ে $P = 6(5) + 6 = 36$ একক। ($6(5)$ বলতে $6\times5$ বোঝায়।)

(খ) $6b + 6 = 48 \Rightarrow 6b = 42 \Rightarrow b = 7$ একক।

লক্ষ করো, একই অক্ষর $b$ দুই জায়গায় দুই ভূমিকায় কাজ করল। $P = 6b+6$ লেখার সময়
$b$ ছিল ফাঁকা ঘর --- যেকোনো ধনাত্মক প্রস্থ বসানো যেত। (খ) অংশে একটা শর্ত চাপিয়ে
দেওয়ার পর $b$ হয়ে গেল অজানা রাশি, আর শর্তটা তাকে একটামাত্র মানে বেঁধে ফেলল।
\end{solbox}

\prob{2} একটি ট্যাক্সিতে প্রথম ২ কিলোমিটার পর্যন্ত ভাড়া ৫০ টাকা নির্দিষ্ট,
তারপর প্রতি কিলোমিটারে ৩০ টাকা। যাত্রার দৈর্ঘ্য $d$ কিলোমিটার ($d \geq 2$) হলে
ভাড়া $F$-কে $d$-এর মাধ্যমে লেখো, এবং $d = 7$ হলে ভাড়া নির্ণয় করো। এই রাশিতে
কোনগুলো ধ্রুবক আর কোনটা চলক?

\begin{sol}
প্রথম $2$ কিলোমিটারের পরে অতিরিক্ত দূরত্ব $d - 2$, তাই
\[
  F = 50 + 30(d - 2) .
\]
$d = 7$ হলে $F = 50 + 30(5) = 200$ টাকা।

\noindent এখানে $50$ ও $30$ ধ্রুবক (ভাড়ার হার বা রেইট (rate), যা যাত্রার সাথে বদলায় না), $d$ চলক, আর $F$
হলো $d$-এর উপর নির্ভরশীল আরেকটি রাশি। এই নির্ভরশীলতাটাই অধ্যায় ৪-এ function
নামে ফিরে আসবে।
\end{sol}

\prob{2} তিনটি সংখ্যা $x_1, x_2, x_3$-এর গড় $12$। এখন প্রতিটি সংখ্যার সাথে $5$ যোগ করে যোগফলকে $2$ দিয়ে গুণ করা হলো। নতুন তিনটি সংখ্যার গড় কত? (উত্তরটা কেবল একটা উদাহরণ দিয়ে বের করো না ; $x_1, x_2, x_3$ ব্যবহার করেই প্রমাণ করো।)

\begin{sol}
নতুন সংখ্যাগুলো $2(x_1 + 5),\ 2(x_2+5),\ 2(x_3+5)$। এদের গড়
\[
  \frac{2(x_1+5) + 2(x_2+5) + 2(x_3+5)}{3}
  = \frac{2\big((x_1+x_2+x_3) + 15\big)}{3}
  = 2\left( \frac{x_1+x_2+x_3}{3} + 5 \right).
\]
ভেতরের ভগ্নাংশটাই পুরোনো গড়, অর্থাৎ $12$। সুতরাং নতুন গড় $= 2(12 + 5) = 34$।

\noindent মনে রেখো, $x_1, x_2, x_3$ এর জন্য তিনটি নির্দিষ্ট সংখ্যা বেছে নিয়ে $34$ পাওয়া কোনো প্রমাণ নয়; সেটা
কেবল একটা পরীক্ষা। প্রমাণ হলো উপরের হিসাব, যেখানে $x_1, x_2, x_3$ শেষ পর্যন্ত
অক্ষর হিসেবেই রয়ে গেছে।
\end{sol}

\section{পদ ও রাশি}

\begin{idea}[সংখ্যার নিয়মগুলোই বীজগণিতের নিয়ম]
আমরা ধরে নিচ্ছি (প্রমাণ করছি না) যে, যেকোনো বাস্তব সংখ্যা $a, b, c$-এর জন্য
\[
  a + b = b + a, \qquad (a+b) + c = a + (b+c), \qquad ab = ba,
\]
\[
  (ab)c = a(bc), \qquad a(b + c) = ab + ac .
\]
এগুলো সংখ্যার জগতের ভিত্তি; এই
সিরিজে আমরা এদের সত্য ধরে নিয়ে শুরু করছি। বীজগণিতের বাকি প্রায় সব কিছু এই পাঁচটা নিয়ম থেকেই বের হয়ে আসে।
\end{idea}

\noindent
কিছু শব্দ ঠিক করে নিই। $7a^2b - 3ab^2 + 4$ রাশিটিতে $+$ ও $-$ চিহ্ন দিয়ে আলাদা হওয়া অংশগুলোকে বলে পদ (term); এখানে পদ তিনটি। প্রতিটি পদে চলকের পূর্বে সংখ্যাসূচক অংশটাকে
বলে সহগ (coefficient) : প্রথম পদের সহগ $7$, দ্বিতীয়টির $-3$। চলকের মাথার উপরে যে ছোট ছোট সংখ্যাগুলো আছে সেগুলো সূচক (exponent) বা ঘাত (power) : $x^2=1\times x \times x, x^3=1\times x\times x\times x, x^1=1\times x, x^0=1$ ইত্যাদি। যে পদে কোনো চলক নেই (অর্থাৎ চলকের ঘাত $0$), তাকে বলে ধ্রুব পদ; এখানে সেটা $4$।

\noindent (উপরের $x$ আর $\times$ গুলো অনেক হিজিবিজি হিজিবিজি লাগছে, তাই না? এই সমস্যা সমাধানে বীজগণিতে গুণন চিহ্ন ($\times$) কে ডট ($\cdot$) লেখা হয়। যেমন $2\times y=2\cdot y, x\times x=x\cdot x$ ইত্যাদি)

\noindent দুটি পদ সদৃশ (like terms) তখনই, যখন তাদের চলক-অংশ হুবহু এক : একই অক্ষর, একই
সূচক। $7a^2b$ ও $-2a^2b$ সদৃশ; $7a^2b$ ও $7ab^2$ নয়।

\begin{idea}[সদৃশ পদ কেন যোগ হয়]
$3x + 5x = 8x$ কোনো আলাদা নিয়ম নয়, এটা বণ্টন বিধিরই প্রয়োগ:
\[
  3x + 5x = (3 + 5)x = 8x .
\]
$3x + 5x^2$ এভাবে জোড়া লাগে না, কারণ $x$ আর $x^2$ এক জিনিস নয়। সাধারণ উৎপাদক হিসেবে $x$ বের করে আনলে পাওয়া যায় $x(3 + 5x)$, যা একটি গুণফল, একক পদ নয়।

তাই ``সদৃশ পদ যোগ করো'' নিয়মটা মুখস্থ করার কিছু নেই। প্রশ্নটা সবসময় একটাই:
সাধারণ উৎপাদক বের করে আনার পর (যদি বের করা যায়) ব্র্যাকেটের ভেতরে কি কেবল সংখ্যা পড়ে থাকে?
\end{idea}

\begin{example}
সরল করো: $\;6a^2b + 2ab^2 - 4a^2b - 5ab^2 + a^2b$। তারপর $a = 3$, $b = -2$
বসিয়ে মান নির্ণয় করো।
\end{example}

\begin{solbox}
$a^2b$ জাতীয় পদগুলো একসাথে, $ab^2$ জাতীয় পদগুলো একসাথে নিই:
\[
  (6 - 4 + 1)a^2b + (2 - 5)ab^2 = 3a^2b - 3ab^2 .
\]
এখন $a = 3$, $b = -2$ বসাই। $a^2b = 9 \cdot (-2) = -18$ এবং
$ab^2 = 3 \cdot 4 = 12$, তাই
\[
  3(-18) - 3(12) = -54 - 36 = -90 .
\]

সরল না করে সরাসরি বসালেও $-90$ আসত। কিন্তু সরল করা রাশিতে পাঁচটার বদলে দুটো পদ,
তাই ভুল করার সুযোগও কম--- সরলীকরণের আসল লাভ এটাই!
\end{solbox}

\begin{remark}
সবচেয়ে বেশি যে ভুলটা হয়: $3x + 5x^2 = 8x^3$। সূচক যোগ হয় \emph{গুণ} করার সময়,
\emph{যোগ} করার সময় নয়। সন্দেহ হলে একটা সংখ্যা বসিয়ে দেখো : $x = 1$ বসালে
বাঁ পাশ $8$, ডান পাশ $8$; কিন্তু $x = 2$ বসালে বাঁ পাশ $26$, ডান পাশ $64$। ভুল প্রমাণ করতে
একটা counterexample যথেষ্ট, আর সেটা খুঁজে পেতে সেকেন্ড দশেক লাগে।
\end{remark}

\prob{2} সরল করো: $\;7a^2b - 3ab^2 + 2a^2b + 5ab^2 - 4a^2b$। তারপর $a = 2$,
$b = -1$ বসিয়ে মান নির্ণয় করো।

\begin{sol}
$a^2b$-এর সহগ $7 + 2 - 4 = 5$, আর $ab^2$-এর সহগ $-3 + 5 = 2$। তাই রাশিটি
$5a^2b + 2ab^2$।

$a = 2$, $b = -1$ বসালে $a^2b = 4(-1) = -4$ এবং $ab^2 = 2(1) = 2$, সুতরাং মান
$5(-4) + 2(2) = -20 + 4 = -16$।
\end{sol}

\prob{2} দেখাও যে $(x+3)(x+5) - (x+4)^2$ রাশিটির মান $x$-এর উপর মোটেও নির্ভর
করে না, এবং সেই মানটি নির্ণয় করো।

\begin{sol}
বণ্টন বিধি দিয়ে দুটো গুণফলই ভেঙ্গে ফেলি:
\[
  (x+3)(x+5) = x^2 + 8x + 15, \qquad (x+4)^2 =(x+4)(x+4) x^2 + 8x + 16 .
\]
বিয়োগ করলে
\[
  (x^2 + 8x + 15) - (x^2 + 8x + 16) = -1 .
\]
$x$-যুক্ত সব পদ কেটে গেছে, তাই যেকোনো $x$-এর জন্য মান $-1$। এটাই পরের অনুচ্ছেদের
বিষয়; এমন একটা সমতা, যা সব $x$-এর জন্য খাটে।
\end{sol}

\section{অভেদ বনাম সমীকরণ}

\begin{idea}[অভেদ সব মানের জন্য, সমীকরণ কেবল কিছু মানের জন্য]
দুটো রাশির মাঝে সমান চিহ্ন বসালে সেটা দুই রকম জিনিস হতে পারে।

\textbf{অভেদ (identity)} : চলকের \emph{সব} গ্রহণযোগ্য মানের জন্য সত্য।
যেমন $(a+b)^2 = a^2 + 2ab + b^2$।

\textbf{সমীকরণ (equation)} : কেবল \emph{কিছু} মানের জন্য সত্য ও বাকিদের জন্য
মিথ্যা, অথবা সকল মানের জন্য সত্য (অর্থাৎ সকল অভেদই একেকটি সমীকরণ)। যেমন $x^2 = 4$, যা কেবল $x = 2$ ও $x = -2$-এ সত্য। আবার $(a+b)^3=a^3+3a^2b+3ab^2+b^3$ একই সাথে একটি অভেদ ও সমীকরণ।

দেখতে দুটোই একরকম। পার্থক্যটা লেখায় নেই, দাবিতে। তবে অনেক সময় অভেদের জন্য $=$ -এর পরিবর্তে $\equiv$ চিহ্ন ব্যবহার করা হয়। 
\end{idea}

\begin{idea}[প্রমাণ ও খণ্ডন এখানে অসমান কাজ]
অধ্যায় ১-এর কথাটা এখানে সরাসরি খাটে।

কোনো একটি নির্দিষ্ট সমতা অভেদ নয়, এটা দেখাতে একটামাত্র counterexample-ই যথেষ্ট।

কিন্তু কোনো একটি সমতা অভেদ, এটা দেখাতে মান বসিয়ে পরীক্ষা করা কখনোই
যথেষ্ট নয়, তা যতগুলো সংখ্যাই বসাও না কেন। কারণ না বসানো মানগুলোর সংখ্যা অসীম, আর সত্য ভাঙার
জন্য তাদের যেকোনো একটাই যথেষ্ট। অভেদ প্রমাণ করতে হয় অক্ষর রেখেই হিসাব করে, অর্থাৎ সমীকরণের এক পাশের রাশিকে বণ্টন বিধি, বিনিময় বিধি, সংযোগ বিধি প্রভৃতির সাহায্যে পরিবর্তন করে করে আরেকপাশের রাশির সমান করার মাধ্যমে।
\end{idea}

\begin{example}
$(x+2)(x+6)$ ও $(x+4)^2$ --- এই দুটি রাশি কি সর্বদা সমান? না হলে, কোন কোন $x$-এর জন্য এদের মান সমান হয়?
\end{example}

\begin{solbox}
দুটোই খুলে ফেলি:
\[
  (x+2)(x+6) = x^2 + 8x + 12, \qquad (x+4)^2 = x^2 + 8x + 16 .
\]
এরা অভেদ নয়; $x = 0$ বসালেই বাঁ পাশ $12$, ডান পাশ $16$। একটা counterexample, এতেই 
কাজ শেষ।

এবার সমান হওয়ার শর্ত খুঁজি:
\[
  x^2 + 8x + 12 = x^2 + 8x + 16 \iff 12 = 16 .
\]
কিন্তু $12 = 16$ মিথ্যা, আর তাতে $x$ নেই-ই। সুতরাং কোনো $x$-এর জন্যই এরা সমান
হয় না --- সমীকরণটির কোনো সমাধান নেই।

তাহলে সমান চিহ্নের তিনটা সম্ভাবনা: সব মানের জন্যই সত্য (অভেদ), কিছু মানের জন্য সত্য (সমাধানসহ
সমীকরণ), কোনো মানের জন্যই সত্য নয় (সমাধানহীন সমীকরণ)।
\end{solbox}

\begin{remark}
স্কুলে ``সূত্র'' নামে যেগুলো মুখস্থ করানো হয়, তার প্রায় সবই আসলে অভেদ;
অর্থাৎ প্রমাণযোগ্য বিবৃতি। $(a+b)^2 = a^2 + 2ab + b^2$ তিন
লাইনে প্রমাণ হয়ে যায়। একবার প্রমাণটা নিজে করে ফেললে ভুলে যাওয়ার ভয়ও থাকে না,
কারণ ভুলে গেলে আবার বের করে নেওয়া যায়। 
\end{remark}

\prob{3} নিচের প্রতিটি সম্পর্কে বলো এটি অভেদ না সমীকরণ। অভেদ হলে প্রমাণ করো;
না হলে একটি counterexample দাও এবং সম্ভব হলে সমাধান নির্ণয় করো।
\begin{enumerate}
  \item $(x-2)(x+2) = x^2 - 4$
  \item $(x+1)^2 = x^2 + 1$
  \item $(a+b)(a-b) + b^2 = a^2$
  \item $x^2 + 1 = 2x$
\end{enumerate}

\begin{sol}
\begin{enumerate}
  \item অভেদ। ভেঙ্গে পাই $x^2 + 2x - 2x - 4 = x^2 - 4$, যা ডান পাশের সমান।
  \item অভেদ নয়। $x = 1$ বসালে বাঁ পাশ $4$, ডান পাশ $2$। সমীকরণ হিসেবে দেখলে
        $x^2 + 2x + 1 = x^2 + 1 \iff 2x = 0 \iff x = 0$; একমাত্র সমাধান $x = 0$।
  \item অভেদ। $(a+b)(a-b) = a^2 - b^2$, তাই বাঁ পাশ $= a^2 - b^2 + b^2 = a^2$।
  \item অভেদ নয়। $x = 0$ বসালে বাঁ পাশ $1$, ডান পাশ $0$। সমীকরণ হিসেবে
        $x^2 - 2x + 1 = 0$, অর্থাৎ $(x-1)^2 = 0$, একমাত্র সমাধান $x = 1$।
\end{enumerate}
লক্ষণীয়, (২) ও (৪) দুটোই অভেদ নয়, কিন্তু দুটোরই সমাধান আছে। ``অভেদ নয়'' আর
``সত্য নয়'' এক কথা নয়।
\end{sol}

\prob{3} প্রমাণ করো যে সব বাস্তব $a, b$-এর জন্য $(a+b)^2 - (a-b)^2 = 4ab$।
তারপর এই অভেদ ব্যবহার করে $103^2 - 97^2$-এর মান কাগজে বড় গুণ না করেই নির্ণয়
করো।

\begin{sol}
দুটো বর্গই খুলি:
\[
  (a+b)^2 = a^2 + 2ab + b^2, \qquad (a-b)^2 = a^2 - 2ab + b^2 .
\]
বিয়োগ করলে $a^2$ ও $b^2$ কেটে যায়:
\[
  (a+b)^2 - (a-b)^2 = 4ab .
\]
এখন $a = 100$, $b = 3$ নিলে $a + b = 103$ এবং $a - b = 97$। সুতরাং
\[
  103^2 - 97^2 = 4(100)(3) = 1200 .
\]
এখানেই অভেদের শক্তি : অক্ষরে প্রমাণ করা একটা কথা; তারপর সুবিধামতো সংখ্যা বসিয়ে
হিসাব সহজ করে ফেলা।
\end{sol}

\section{সূচক}

\begin{idea}[সূচক]
ধনাত্মক পূর্ণসংখ্যা $n$-এর জন্য $a^n$ মানে $n$ সংখ্যক $a$-এর গুণফল:
$a^n = a \cdot a \cdots a$। এই একটামাত্র সংজ্ঞা থেকেই নিচের নিয়মগুলো বেরিয়ে আসে; প্রতিটিই আসলে উৎপাদক গোনা।

\medskip
$a^m \cdot a^n = a^{m+n}$: বাঁ পাশে প্রথমে $m$টি, তারপর $n$টি উৎপাদক, মোট
$m+n$টি।

$(a^m)^n = a^{mn}$: $m$টি উৎপাদকের গুচ্ছ $n$ বার, মোট $mn$টি।

$(ab)^n = a^n b^n$: $n$টি $ab$ পাশাপাশি লিখে ক্রম বিনিময় করে সব $a$ একদিকে ও
সব $b$ একদিকে সরিয়ে নিলেই হলো।

$\dfrac{a^m}{a^n} = a^{m-n}$ (যেখানে $m > n$, $a \neq 0$): উপরে-নিচে $n$টি করে
উৎপাদক কেটে গিয়ে উপরে $m - n$টি পড়ে থাকে।
\end{idea}

\begin{idea}[$a^0$ ও ঋণাত্মক সূচক আবিষ্কার নয়, সিদ্ধান্ত]
$a^0$ মানে ``শূন্যবার $a$ গুণ'' --- কথাটার নিজস্ব কোনো অর্থ নেই। $a^{-3}$-এরও নেই।
তাহলে এদের মান কীভাবে ঠিক হলো?

আমরা \emph{বেছে} নিয়েছি, এবং শর্ত ছিল একটাই: প্রথম নিয়মটা যেন ভেঙে না পড়ে।

$a^m \cdot a^0 = a^{m+0} = a^m$ সত্য রাখতে চাইলে, $a \neq 0$ হলে $a^m \neq 0$,
তাই দুই পাশকে $a^m$ দিয়ে ভাগ করে পাই $a^0 = 1$। অন্য কোনো মান নেওয়ার উপায় নেই।

একইভাবে $a^n \cdot a^{-n} = a^0 = 1$ রাখতে চাইলে
\[
  a^{-n} = \frac{1}{a^{n}} \qquad (a \neq 0).
\]
এই দুটো সংজ্ঞা মেনে নিলে উপরের চারটি নিয়মই \emph{সব} পূর্ণসংখ্যা $m, n$-এর জন্য
খাটে, শুধু $a \neq 0$ শর্তে। গণিতে এভাবে সংজ্ঞা বাড়ানো খুব সাধারণ ঘটনা: পুরোনো
নিয়ম টিকিয়ে রাখাই নতুন সংজ্ঞা বাছার মাপকাঠি।
\end{idea}

\begin{example}
সরল করো এবং উত্তর কেবল ধনাত্মক সূচকে লেখো:
$\;\dfrac{(2x^{-2}y^{3})^{3}}{4x^{-5}y^{4}}$, যেখানে $x \neq 0$, $y \neq 0$।
\end{example}

\begin{solbox}
লব থেকে শুরু করি:
\[
  (2x^{-2}y^{3})^{3} = 2^{3} \, x^{-6} \, y^{9} = 8x^{-6}y^{9} .
\]
এখন ভাগ করি, প্রতিটি অক্ষরের সূচক বিয়োগ করে:
\[
  \frac{8x^{-6}y^{9}}{4x^{-5}y^{4}}
  = 2 \, x^{-6-(-5)} \, y^{9-4}
  = 2x^{-1}y^{5}
  = \frac{2y^{5}}{x} .
\]
মাঝের ধাপে $-6 - (-5) = -1$; ঋণাত্মক সূচক বিয়োগ করার সময় চিহ্নেই সবচেয়ে বেশি
ভুল হয়।
\end{solbox}

\begin{remark}
দুটো ফাঁদ। প্রথমত, $-2^{4}$ আর $(-2)^{4}$ এক নয়: সূচক কেবল তার ঠিক আগের রাশিটার
উপর কাজ করে, তাই $-2^{4} = -16$ কিন্তু $(-2)^{4} = 16$। দ্বিতীয়ত, $a^{m}b^{n}\neq(ab)^{m+n}$, কারণ সূচকের নিয়মগুলো একই ভিত্তির (base) জন্য, ভিন্ন ভিত্তির জন্য নয়।

ভগ্নাংশ সূচক, যেমন $a^{1/2}$, এখানে ইচ্ছে করেই বাদ রাখা হয়েছে। ওটার জন্য আগে
function ও inverse দরকার, তাই ভগ্নাংশ সূচক ও logarithm আসবে অধ্যায় ৮-এ।
\end{remark}

\prob{2} সরল করো, উত্তর কেবল ধনাত্মক সূচকে লেখো: $\;(3x^{2}y^{-3})(2x^{-5}y^{4})$,
যেখানে $x \neq 0$, $y \neq 0$।

\begin{sol}
সহগ আলাদা, প্রতিটি অক্ষর আলাদা:
\[
  (3 \cdot 2)\, x^{2 + (-5)} \, y^{-3 + 4} = 6x^{-3}y = \frac{6y}{x^{3}} .
\]
\end{sol}

\prob{2} $n$ যেকোনো পূর্ণসংখ্যা হলে দেখাও যে
$\;\dfrac{2^{\,n+2} - 2^{\,n}}{2^{\,n+1}} = \dfrac{3}{2}$, অর্থাৎ ভগ্নাংশটির মান
$n$-এর উপর নির্ভর করে না।

\begin{sol}
লব থেকে $2^{n}$ সাধারণ উৎপাদক হিসেবে বের করে আনি:
\[
  2^{\,n+2} - 2^{\,n} = 2^{n} \cdot 2^{2} - 2^{n} = 2^{n}(4 - 1) = 3 \cdot 2^{n} .
\]
হর $2^{\,n+1} = 2 \cdot 2^{n}$। সুতরাং
\[
  \frac{3 \cdot 2^{n}}{2 \cdot 2^{n}} = \frac{3}{2},
\]
কারণ $2^{n} \neq 0$, তাই সেটা কাটা বৈধ। $n$ ধনাত্মক না ঋণাত্মক, তাতে কিছু আসে
যায় না; নিয়মগুলো সব পূর্ণসংখ্যার জন্যই খাটে।
\end{sol}

\prob{3} দেখাও যে সূচকের প্রথম নিয়মটি ($a^{m}a^{n} = a^{m+n}$) $n = 0$-এর জন্যও
সত্য রাখতে চাইলে $a \neq 0$ হলে $a^{0} = 1$ ছাড়া অন্য কোনো সংজ্ঞা সম্ভব নয়।
তারপর ব্যাখ্যা করো, $a = 0$ হলে এই যুক্তিটা ঠিক কোথায় ভেঙে পড়ে।

\begin{sol}
নিয়মটি $n = 0$-এ প্রয়োগ করলে
\[
  a^{m} \cdot a^{0} = a^{m + 0} = a^{m} .
\]
$a \neq 0$ হলে $a^{m} \neq 0$, তাই দুই পাশকে $a^{m}$ দিয়ে ভাগ করা যায় এবং পাই
$a^{0} = 1$। অর্থাৎ এই সংজ্ঞাটা আমাদের হাতে ছিল না, নিয়মটাই তাকে বেঁধে দিয়েছে।

\noindent $a = 0$ হলে $a^{m} = 0$ (ধনাত্মক $m$-এর জন্য), আর সমীকরণটা দাঁড়ায়
$0 \cdot 0^{0} = 0$। এটা $0^{0}$-এর \emph{যেকোনো} মানের জন্যই সত্য, তাই কোনো
একটা নির্দিষ্ট মান বেছে নেওয়ার যুক্তি পাওয়া যায় না। শূন্য দিয়ে ভাগ করা যায় না বলেই
ধাপটা আটকে যায়, আর এ কারণেই $0^{0}$-কে অসংজ্ঞায়িত রাখা হয়।
\end{sol}

\section{সমীকরণ সমাধান}

\begin{idea}[দুই পাশে একই কাজ করা]
একটি সমীকরণ সমাধান করা মানে তাকে ধাপে ধাপে সহজ করা, এমনভাবে যেন প্রতিটি ধাপে
সমীকরণের দুইপাশ একই থাকে। যে ধাপগুলো এটা নিশ্চিত করে:

\medskip
দুই পাশে একই সংখ্যা যোগ বা বিয়োগ করা;

দুই পাশকে একই \emph{অশূন্য} সংখ্যা দিয়ে গুণ বা ভাগ করা।

\medskip
এদের প্রতিটিই ফিরিয়ে আনা যায়, অনেকটা Undo করার মতো : যোগ করলে বিয়োগ করে, গুণ করলে ভাগ করে। তাই
পুরোনো সমীকরণ ও নতুন সমীকরণ পরস্পরের সমতুল্য থাকে, অর্থাৎ এদের মাঝে $\Leftrightarrow$
বসে।

শূন্য দিয়ে গুণ করলে সমীকরণটিকে আবার ফিরিয়ে আনা যায় না, তাই সেটা নিষিদ্ধ।
\end{idea}

\begin{idea}[রৈখিক সমীকরণ]
$ax + b = 0$ আকারের সমীকরণকে বলে এক চলকের রৈখিক (linear) সমীকরণ। $a \neq 0$ হলে
\[
  ax + b = 0 \iff ax = -b \iff x = -\frac{b}{a},
\]
অর্থাৎ ঠিক একটি সমাধান। কিন্তু $a = 0$ হলে সমীকরণটা দাঁড়ায় $b = 0$, যেখানে $x$
নেই-ই। তখন দুটো সম্ভাবনা: $b \neq 0$ হলে কোনো সমাধান নেই, আর $b = 0$ হলে
সব সংখ্যাই $x$-এর সমাধান; অর্থাৎ ওটা আসলে সমীকরণ নয়, অভেদ।

লক্ষ করো, ``$a$ দিয়ে ভাগ করো'' বলার আগে $a \neq 0$ যাচাই করতেই হয়। এই ছোট
শর্তটা না দেখে এড়িয়ে যাওয়া পরে অনেক বড় ভুলের উৎস।
\end{idea}

\begin{remark}
সব ধাপ ফিরিয়ে আনা যায় না, আর তখনই ভুতুড়ে সমাধান জন্ম নেয়। যেমন
$\dfrac{x^{2}-4}{x-2} = 4$ সমীকরণে দুই পাশকে $x - 2$ দিয়ে গুণ করলে পাই
$x^{2} - 4 = 4x - 8$, অর্থাৎ $x^{2} - 4x + 4 = 0$, অর্থাৎ $x = 2$। কিন্তু মূল
সমীকরণে $x = 2$ বসানোই যায় না; হর শূন্য হয়ে যায়। আসলে $x = 2$ হলে আমরা শূন্য
দিয়ে গুণ করেছিলাম, আর সেই ধাপটা ফেরানো যায় না। মূল সমীকরণের কোনো সমাধান নেই।

নিয়ম: চলকযুক্ত কোনো রাশি দিয়ে গুণ করলে শেষে প্রতিটি উত্তর মূল সমীকরণে বসিয়ে
যাচাই করতেই হবে।
\end{remark}

\begin{example}
সমাধান করো: $\;\dfrac{3x - 1}{4} - \dfrac{x + 2}{3} = \dfrac{1}{2}$।
\end{example}

\begin{solbox}
হরগুলোর ল.সা.গু. $12$। দুই পাশকে $12$ দিয়ে গুণ করি --- $12 \neq 0$, তাই এটা
বৈধ ও ফিরিয়ে আনা যায়:
\[
  3(3x - 1) - 4(x + 2) = 6 .
\]
বাঁ পাশ খুলি: $9x - 3 - 4x - 8 = 5x - 11$। সুতরাং
\[
  5x - 11 = 6 \iff 5x = 17 \iff x = \frac{17}{5} .
\]
যাচাই: $\dfrac{3(17/5) - 1}{4} = \dfrac{46/5}{4} = \dfrac{23}{10}$ এবং
$\dfrac{17/5 + 2}{3} = \dfrac{27/5}{3} = \dfrac{9}{5}$। বিয়োগ করলে
$\dfrac{23}{10} - \dfrac{18}{10} = \dfrac{5}{10} = \dfrac{1}{2}$। মিলে গেছে।
\end{solbox}

\subsection*{দুই চলকের সরল সহসমীকরণ}

\noindent
এবার দুটি অজানা রাশি। সাধারণ রূপ
\[
  a_1 x + b_1 y + c_1 = 0, \qquad a_2 x + b_2 y + c_2 = 0 .
\]
একটিমাত্র সমীকরণে দুটি অজানা থাকলে সমাধান অসংখ্য; দুটি সমীকরণ একসাথে থাকলে তবেই
সাধারণত একটিমাত্র জোড়া $(x, y)$ পাওয়া যায়। মূল কৌশল সবসময় একটাই : একটি চলককে
সরিয়ে ফেলে সমস্যাটাকে এক চলকে নামিয়ে আনা, যেটা সমাধান করা আমরা এইমাত্র শিখলাম।

\begin{idea}[তিনটি পদ্ধতি, একই ভাবনা]
\textbf{প্রতিস্থাপন (substitution)} --- একটি সমীকরণ থেকে একটি চলককে অন্যটির
মাধ্যমে প্রকাশ করে সেটা অন্য সমীকরণে বসিয়ে দাও।

\textbf{অপনয়ন (elimination)} --- দুটি সমীকরণকে উপযুক্ত অশূন্য সংখ্যা দিয়ে গুণ
করে এমন করো যেন কোনো এক চলকের সহগ দুই জায়গায় সমান হয়; তারপর বিয়োগ করে সেই চলক
মুছে ফেলো।

\textbf{আড়গুণন (cross-multiplication)} --- অপনয়ন একবার সাধারণ অক্ষরে করে রেখে
পাওয়া সূত্র, যা প্রতিবার নতুন করে না করে সরাসরি ব্যবহার করা যায়।
\end{idea}

\noindent
আড়গুণনের সূত্রটা আকাশ থেকে পড়ে না; এটা অপনয়ন, শুধু সংখ্যার বদলে অক্ষর নিয়ে।
সমীকরণ দুটো লিখি $a_1x + b_1y = -c_1$ ও $a_2x + b_2y = -c_2$। প্রথমটিকে $b_2$
আর দ্বিতীয়টিকে $b_1$ দিয়ে গুণ করে বিয়োগ করলে $y$ মুছে যায়:
\[
  (a_1b_2 - a_2b_1)\,x = b_1c_2 - b_2c_1 .
\]
আবার প্রথমটিকে $a_2$ আর দ্বিতীয়টিকে $a_1$ দিয়ে গুণ করে বিয়োগ করলে $x$ মুছে যায়:
\[
  (a_1b_2 - a_2b_1)\,y = c_1a_2 - c_2a_1 .
\]
এখন যদি $a_1b_2 - a_2b_1 \neq 0$ হয়, তবেই কেবল ভাগ করা যায়, এবং তিনটে ফল একসাথে
লেখা যায়:
\[
  \frac{x}{\,b_1c_2 - b_2c_1\,} \;=\; \frac{y}{\,c_1a_2 - c_2a_1\,}
  \;=\; \frac{1}{\,a_1b_2 - a_2b_1\,} .
\]

\begin{idea}[কখন সমাধান আছে]
সব কিছু নির্ভর করছে $D = a_1b_2 - a_2b_1$ রাশিটির উপর।

$D \neq 0$ হলে ঠিক একটি সমাধান, আর সেটা উপরের সূত্র দিচ্ছে।

$D = 0$ হলে ভাগ করা যায় না, এবং তখন দুটো সম্ভাবনা :  হয় সমীকরণ দুটি পরস্পরের
গুণিতক, তখন একটি আসলে অন্যটিরই পুনরুক্তি এবং সমাধান অসংখ্য; নয়তো তারা পরস্পরবিরোধী
এবং কোনো সমাধান নেই।
\end{idea}

\begin{example}
সমাধান করো: $\;2x + 5y = 16$ এবং $3x - 2y = 5$। প্রতিস্থাপন ও অপনয়ন --- দুই
পদ্ধতিতেই করো।
\end{example}

\begin{solbox}
\textbf{অপনয়ন।} $y$ মোছার জন্য প্রথমটিকে $2$ আর দ্বিতীয়টিকে $5$ দিয়ে গুণ করি:
\[
  4x + 10y = 32, \qquad 15x - 10y = 25 .
\]
যোগ করলে $19x = 57$, তাই $x = 3$। প্রথম সমীকরণে বসিয়ে $6 + 5y = 16$, অর্থাৎ
$y = 2$।

\medskip
\textbf{প্রতিস্থাপন।} দ্বিতীয় সমীকরণ থেকে $3x = 5 + 2y$, অর্থাৎ
$x = \dfrac{5 + 2y}{3}$। এটা প্রথম সমীকরণে বসাই:
\[
  2 \cdot \frac{5 + 2y}{3} + 5y = 16 .
\]
দুই পাশকে $3$ দিয়ে গুণ করে $10 + 4y + 15y = 48$, অর্থাৎ $19y = 38$, অর্থাৎ
$y = 2$। তারপর $x = \dfrac{5 + 4}{3} = 3$।

\medskip
যাচাই: $2(3) + 5(2) = 16$ এবং $3(3) - 2(2) = 5$। দুটোই মিলেছে।

দুই পদ্ধতিতে একই উত্তর আসাটা কাকতাল নয় --- দুটোই কেবল সমতুল্য রূপান্তর ব্যবহার
করেছে, তাই সমাধানের তালিকা বদলাতেই পারে না।
\end{solbox}

\prob{2} সমাধান করো: $\;\dfrac{2x - 1}{3} - \dfrac{x + 2}{4} = \dfrac{1}{6}$।
উত্তরটি মূল সমীকরণে বসিয়ে যাচাই করো।

\begin{sol}
হরগুলোর ল.সা.গু. $12$; দুই পাশকে $12$ দিয়ে গুণ করি:
\[
  4(2x - 1) - 3(x + 2) = 2 .
\]
খুলে পাই $8x - 4 - 3x - 6 = 2$, অর্থাৎ $5x - 10 = 2$, অর্থাৎ $5x = 12$, অর্থাৎ
$x = \dfrac{12}{5}$।

যাচাই: $\dfrac{2(12/5) - 1}{3} = \dfrac{19/5}{3} = \dfrac{19}{15}$ এবং
$\dfrac{12/5 + 2}{4} = \dfrac{22/5}{4} = \dfrac{11}{10}$। বিয়োগ করলে
$\dfrac{38}{30} - \dfrac{33}{30} = \dfrac{5}{30} = \dfrac{1}{6}$।
\end{sol}

\prob{3} সমাধান করো: $\;3x + 4y = 18$ এবং $5x - 2y = 4$। প্রথমে অপনয়ন দিয়ে
সমাধান করো, তারপর আড়গুণনের সূত্র দিয়ে একই উত্তর পাও কি না দেখো।

\begin{sol}
\textbf{অপনয়ন।} দ্বিতীয় সমীকরণকে $2$ দিয়ে গুণ করলে $10x - 4y = 8$। প্রথমটির
সাথে যোগ করলে $13x = 26$, তাই $x = 2$। তারপর $3(2) + 4y = 18$ থেকে $4y = 12$,
অর্থাৎ $y = 3$।

\medskip
\textbf{আড়গুণন।} সমীকরণ দুটোকে সূত্রের আকারে লিখি:
$3x + 4y - 18 = 0$ ও $5x - 2y - 4 = 0$। তাহলে $a_1 = 3$, $b_1 = 4$, $c_1 = -18$,
$a_2 = 5$, $b_2 = -2$, $c_2 = -4$। হিসাব করি:
\[
  b_1c_2 - b_2c_1 = (4)(-4) - (-2)(-18) = -16 - 36 = -52,
\]
\[
  c_1a_2 - c_2a_1 = (-18)(5) - (-4)(3) = -90 + 12 = -78,
\]
\[
  a_1b_2 - a_2b_1 = (3)(-2) - (5)(4) = -6 - 20 = -26 .
\]
শেষেরটি শূন্য নয়, তাই সূত্র খাটে:
\[
  x = \frac{-52}{-26} = 2, \qquad y = \frac{-78}{-26} = 3 .
\]
যাচাই: $5(2) - 2(3) = 4$। মিলেছে।
\end{sol}

\prob{3} $k$-এর কোন মানের জন্য $\;2x + 3y = 7$ এবং $4x + ky = 14$ ---
সমীকরণজোটটির অসংখ্য সমাধান থাকবে? $k$-এর এমন কোনো মান কি আছে যার জন্য জোটটির
কোনো সমাধান থাকবে না? উত্তরের কারণ দেখাও।

\begin{sol}
এখানে $a_1 = 2$, $b_1 = 3$, $a_2 = 4$, $b_2 = k$, তাই
\[
  D = a_1b_2 - a_2b_1 = 2k - 12 .
\]
$D \neq 0$, অর্থাৎ $k \neq 6$ হলে ঠিক একটি সমাধান।

$k = 6$ হলে দ্বিতীয় সমীকরণ দাঁড়ায় $4x + 6y = 14$, যা প্রথমটিকে হুবহু $2$ দিয়ে
গুণ করলেই পাওয়া যায়। অর্থাৎ দ্বিতীয় সমীকরণ নতুন কোনো তথ্য দিচ্ছে না; কার্যত
একটিমাত্র সমীকরণে দুটি অজানা, তাই সমাধান অসংখ্য। সুতরাং $k = 6$।

সমাধানহীন হওয়ার জন্য দরকার $D = 0$, অথচ সমীকরণ দুটি পরস্পরবিরোধী। কিন্তু $D = 0$
হয় কেবল $k = 6$-এ, আর সেখানে আমরা দেখলাম তারা পরস্পরবিরোধী নয় --- একেবারে
সমতুল্য। তাই $k$-এর কোনো মানের জন্যই জোটটি সমাধানহীন হয় না।

ডান পাশের $14$ যদি অন্য কিছু হতো, যেমন $15$, তবে $k = 6$-এ জোটটি সমাধানহীন হয়ে
যেত। শুধু সহগের দিকে তাকিয়ে সিদ্ধান্ত নেওয়া যায় না, ধ্রুব পদও গোনায় ধরতে হয়।
\end{sol}

\begin{remark}
$D = a_1b_2 - a_2b_1$ রাশিটা এখানে হঠাৎ গজিয়ে ওঠা একটা হিসাব মনে হচ্ছে, কিন্তু
এটা সিরিজে বারবার ফিরে আসবে। অধ্যায় ৬-এ দেখবে প্রতিটি সমীকরণ আসলে একটি
সরলরেখা, আর $D = 0$ মানে রেখা দুটি সমান্তরাল. তাই হয় তারা মেলে না (সমাধান
নেই), নয়তো তারা একই রেখা (অসংখ্য সমাধান)। অধ্যায় 20-এ একই রাশি determinant
নামে ফিরবে, এবং তখন দেখবে ওটাই দুটি vector দিয়ে তৈরি সামান্তরিকের ক্ষেত্রফল :
শূন্য হওয়া মানে vector দুটি একই রেখা বরাবর।
\end{remark}

\end{document}
